\section[The genuine directed IH weak-C frontier]{The genuine directed IH weak-$C$ frontier}
\label{sec:c2-weak-ih-frontier}

This section gives a source reduction, records failure of the proposed
four-product method, and applies the reviewed rational-curve admission
to the displayed families. It gives no generic multiparameter
admission. Actual source angles and independent testing variables
remain distinct.

Write $z=u_c^2$, $y=v_c^2$, and use the exact contact quantities
\[
 H=v_cY_c=zh,\quad L=v_c(u_cW_c-Y_c)=yt,\quad
 B=H+L=u_cv_cW_c=zb,\quad R=-u_cX_c=zr.
\]
Thus $0<H<1$, $0<R<z$, $L>0$, and
\[
 L<y(1+R),\quad B>yR,\quad
 G=B+R-z(1-H)>0.
\]
Set $S=1+R+L$ and
\begin{align*}
 \widetilde\Delta={}&(1+z)B^2+2(1+z)BR+(1+z)(1+y)B\\
 &+(z-y)R^2+2z(1+y)R+z(1+y),\\
 \widehat N={}&(1+z)(B+R)(H+y)-(1+y)(z-R)L.
\end{align*}
These are exactly $z\Delta,z^2N_{\rm src}$. In particular,
\[
 \widehat N=(1+z)(B+R)H+yG+z\{y(1+R)-L\}+RL(1+y)>0.
\]
The original coefficients, without an approximation, are
\begin{align*}
 C&=\frac{(1+R)\widehat N}{2L\widetilde\Delta},&
 \alpha&=\frac{(B+R)\widehat N}{2(1-H)\widetilde\Delta},\\
 A&=\frac{y(1+R)(B+R)S}{2H\widetilde\Delta},&
 D&=\frac{(1+z)(B+R)(B-yR)S}
 {2\{y(1+R)-L\}\widetilde\Delta},\\
 m&=\frac{z(1+y)(z-R)(1+R)S}{2G\widetilde\Delta},&
 n+\frac12&=\frac{z(B+R)(1+R)S}{2R\widetilde\Delta}.
                                                        \tag{WI1}
\end{align*}
The checker verifies all six against the unchanged original source.

\paragraph{Uniform source reduction for IH.}
Suppose $z$ is positive infinitesimal, $y$ is positive unbounded,
and $n\ge0$. Then $C$ is weak if and only if $B$ is unbounded.
To prove the nontrivial direction, suppose $B$ is bounded. Since
$R<z$ and $G>0$, one has $z\le4(B+R)$: use $B\ge1/2$
or, otherwise, $B+R>z(1-H)>z/2$. Direct division in (WI1)
gives bounded $\widetilde\Delta/[y(B+R)]$. Moreover,
\[
 \frac{\widehat N}{y(B+R)}
 =(1+z)(1+H/y)
 -(1+1/y)\frac{(z-R)L}{B+R}
\]
has residue $1$, since $L\le B$ and $z$ is infinitesimal.
Consequently $C\ge c/L$ for a positive real $c$, which is strong
because $L$ is bounded.

Now let $B$ be unbounded. Put
\[
 \rho=B/y,\qquad e=1+R-L/y>0.
\]
Here $0<\rho<1+R+H/y$, so $\rho$ is bounded and its residue
lies in $[0,1]$. The following exact normalized quantities have
residue $1$:
\[
 \frac{\widetilde\Delta}{B(B+y)},\qquad
 \frac{\widehat N}{yB},\qquad
 \frac{L}{B},\qquad \frac{S}{B},\qquad\frac{G}{B}.
                                                        \tag{WI2}
\]
They give
\[
 n+\tfrac12=\frac{\rho}{2r(1+\rho)}(1+o(1)).
\]
The inequality $n\ge0$ therefore gives $r/\rho$ bounded and
$r\le\rho/[1+\rho](1+o(1))\le1/2+o(1)$.
In particular $1-r$ is a positive unit, and
$R/\rho=z(r/\rho)$ is infinitesimal. The normalized leading
relations, uniform at arbitrary openings $H,1-H,e$, are
\begin{align*}
 C&=\frac{1+o(1)}{2B(1+\rho)},&
 A&=\frac{B(1+o(1))}{2H(1+\rho)},\\
 \alpha&=\frac{B(1+o(1))}{2(1-H)(1+\rho)},&
 m&=\frac{z^2(1-r)(1+o(1))}{2B(1+\rho)},\\
 D&=\frac{\rho^2(1+o(1))}{2e(1+\rho)}.
                                                        \tag{WI3}
\end{align*}
Here each $1+o(1)$ means a positive coefficient unit of residue one;
it never denotes a comparison with a testing infinitesimal. In
particular $A,\alpha$ are unbounded, $m,C$ are infinitesimal, and
$Cm$ is infinitesimal for every weak-$C$ IH source. Thus the four
reciprocal-product exit can never admit this entire directed leaf.
If $\rho$ is infinitesimal then $e$ is a unit and $D$ is weak;
if $\rho$ is a unit then $D$ is strong. Literal $n=0$ remains
allowed throughout the reduction. No lower bound on $r$ was used.

\paragraph{A physical family violating three of the four products.}
For positive infinitesimal $\varepsilon$, choose
\[
 z=\varepsilon^2,\quad y=\varepsilon^{-4},\quad
 H=\tfrac12,\quad B=\varepsilon^{-1},\quad r=\varepsilon^3/4,
 \qquad L=B-H,\quad R=zr.
\]
All exact positive-source gaps hold. The unchanged six coefficients
satisfy
\[
 n\to\tfrac32,\quad m\sim\tfrac12\varepsilon^5,\quad
 C\sim\tfrac12\varepsilon,\quad A\sim\varepsilon^{-1},\quad
 D\sim\tfrac12\varepsilon^6,\quad\alpha\sim\varepsilon^{-1}.
\]
Thus $Cn,Cm,CD$ are infinitesimal, while $CA\to1/2$.
The contact is exactly
\[
 (u_c,v_c;W_c,X_c,Y_c,Z_c)
 =(\varepsilon,\varepsilon^{-2};1,-\varepsilon^4/4,
   \varepsilon^2/2,1).
\]
This is physical, not merely a positive stationary tuple. We use the
inherited stationary-contact sufficiency argument from the general
original-parent geometry in Section~\ref{sec:physical}:
at a finite nonzero exact source contact the signed core restricted to
$Z=0$ is positive definite, so the state matrix has three positive
transverse eigenvalues. Kernel and angular stationarity imply a
stationary determinant zero; its nonnegative Laurent outer
coefficients make that zero a global minimum. The inherited determinant
criterion and homogeneous endpoint continuity prove physicality.
This is a first-order real algebraic implication and transfers to the
coefficient real closed field. No tensor-SOS assertion is imported
from this physicality argument.


\paragraph{The displayed rational curves are individually admitted.}
The exact $\varepsilon$-family above, the reflected HH family below,
and the radial-obstruction family with $r=\varepsilon^6/4$ all have
coefficients in $\mathbb R(\varepsilon)$. The checked positive leading
terms of every strict source gap and coefficient imply one common
real interval $(0,b)$ on which all these signs hold. The inherited
stationary-contact argument then proves physicality throughout that
interval. Their denominators have no zero after $b$ is decreased if
necessary. A positive polynomial common denominator, obtained by
multiplying the squared denominator factors including the powers of
$\varepsilon$, clears all poles. The cleared full parent is
nonnegative on the closed interval by continuity. Therefore
Corollary~\ref{cor:rational-physical-source-curve} from
\path{rational_curve_tensor_sos.tex} gives tensor SOS at every
independent projective test and every real or complex state for each
of these exact curves, including positive infinitesimal
$\varepsilon\in K_c$.

Consequently these curves demonstrate failure of the proposed
four-product or initial toric method, not nonemptiness of a remaining
source leaf after all admissions. They are individually all-tests
admitted. The generic multiparameter IH and HH domains are not closed
by this curve application.
